% --- Latar Belakang ---
\section{Latar Belakang}
Industri \textit{Fast Moving Consumer Goods} (FMCG) merupakan sektor dengan pertumbuhan paling pesat secara global. Pasar FMCG dunia bernilai USD 14,1 triliun pada tahun 2024 dengan CAGR 5,4\% hingga tahun 2034 \autocite{researchandmarkets2025}, sementara kawasan Asia-Pasifik menyumbang lebih dari 40\% total penjualan global. Di Indonesia, nilai pasar FMCG melampaui USD 100 miliar dengan CAGR 7,6\% (2021--2025), khususnya pada segmen kecantikan dan perawatan diri yang mencatat penjualan senilai Rp36 triliun sepanjang Januari hingga Juli 2025 \autocite{kadin2025}. Dinamika ini mendorong kebutuhan sistem pergudangan yang semakin efisien, seiring pasar \textit{Smart Warehouse} global yang tumbuh dari USD 28 miliar (2024) menuju USD 90,21 miliar (2033) pada CAGR 14,2\%, dengan Asia-Pasifik memimpin laju adopsi di CAGR 18,4\% \autocite{precedenceresearch2025}. Namun, sistem pergudangan konvensional pada industri FMCG masih mengandung inefisiensi signifikan: 62\% keterlambatan pengiriman bersumber dari aktivitas \textit{order picking} yang tidak terkoordinasi \autocite{framinan2025}, dengan kekurangan SDM dalam proses \textit{picking} mencatat nilai keparahan (\textit{mean severity}) 3,79 dari skala 5 \autocite{habazin2017}. Kondisi ini mempertegas urgensi transformasi digital dalam manajemen pergudangan modern.

PT Paragon Technology and Innovation (Paragon Corp) merupakan salah satu perusahaan FMCG terkemuka di Indonesia yang bergerak di bidang manufaktur dan distribusi produk kecantikan serta perawatan diri. Perusahaan ini mengoperasikan lebih dari 40 \textit{Distribution Center} (DC) yang tersebar di Indonesia dan Malaysia \autocite{undip2019}, dengan tiga jenis fasilitas gudang utama, yaitu gudang bahan baku, gudang \textit{packaging}, dan gudang barang jadi (\textit{Finished Goods}). Gudang barang jadi menjadi titik kritis dalam rantai distribusi karena menjadi tahap akhir sebelum produk dikirimkan kepada mitra ritel dan konsumen. Guna merespons kebutuhan operasional yang terus meningkat sekaligus selaras dengan agenda transformasi digital dalam konsep \textit{Society} 5.0, Paragon Corp menginisiasi proyek perancangan sistem \textit{Smart Warehouse} yang mencakup integrasi perangkat keras robotika dan layanan perangkat lunak secara menyeluruh pada gudang barang jadinya \autocite{docB100}.

Berdasarkan analisis proses bisnis pada gudang barang jadi Paragon Corp, ditemukan permasalahan kritis pada tahap \textit{stocking} dan \textit{picking} yang masih bergantung sepenuhnya pada operator manusia. Kondisi ini menimbulkan \textit{bottleneck} operasional dengan tingkat \textit{overcapacity} di atas 80\% disertai \textit{waiting time} yang tinggi. Metode \textit{picking} manual mencatat tingkat kesalahan 1--3\%, angka yang 8--10 kali lebih tinggi dibanding lingkungan terotomasi \autocite{netsuite2024}, dengan biaya per insiden berkisar USD 85--220 \autocite{scmchamps2024}. Dampaknya meluas: 65\% pelanggan berhenti berbelanja setelah 2--3 keterlambatan pengiriman \autocite{conexiom2024}, sementara distorsi inventori global akibat ketidakakuratan data diestimasi merugikan USD 1,77 triliun pada tahun 2023 \autocite{ihlgroup2023}. \textcite{vangeest2021} menegaskan bahwa transisi ke \textit{Smart Warehouse}---melalui pilar \textit{Automation and Robotics}, \textit{IoT-based Visibility}, \textit{AI-driven Optimization}, dan \textit{Cyber-Physical Integration}---mampu meningkatkan efisiensi operasional gudang hingga 35\% dibanding sistem konvensional.

Dalam operasionalnya saat ini, Paragon Corp memanfaatkan dua platform teknologi utama untuk pencatatan dan pelacakan barang di gudang. Platform pertama adalah sistem \textit{Enterprise Resource Planning} (ERP) Odoo yang berfungsi sebagai pusat pengelolaan data transaksi bisnis dan inventori secara korporat. Platform kedua adalah \textit{Warehouse Management System} (WMS) yang digunakan untuk mendukung manajemen aktivitas operasional di tingkat gudang secara lebih spesifik. Selain kedua platform tersebut, Paragon Corp juga telah mengadopsi robot \textit{Automated Guided Vehicle} (AGV) untuk sebagian proses mobilisasi barang jadi dari area produksi menuju \textit{Loading Dock}. Sistem kerja keseluruhan menerapkan pendekatan \textit{Pull System}, yakni penggerakan barang berdasarkan permintaan aktual guna meminimalkan penumpukan stok dan mengoptimalkan responsivitas rantai pasok.

Meskipun kedua platform tersebut telah berperan dalam mendukung operasional, integrasi antara ERP Odoo dan WMS masih bersifat semi-manual, sehingga setiap proses sinkronisasi data memerlukan intervensi langsung dari operator gudang. Ketiadaan integrasi \textit{real-time} antara kedua sistem ini mengakibatkan seringnya terjadi \textit{mismatch} data stok antara kondisi fisik di lapangan dengan catatan digital pada sistem pusat. Hal ini diperparah oleh absennya mekanisme pencatatan barang yang bersifat atomik dan otomatis, sehingga area penempatan barang sementara kerap tidak tercatat secara sistematis dalam WMS. Akibatnya, muncul permasalahan \textit{mistracking} barang yang membuat staf gudang tidak dapat mengetahui posisi aktual suatu barang secara akurat. Selain itu, implementasi robot AGV yang masih terbatas pada area produksi menyebabkan proses \textit{inbound} dan perpindahan barang di area gudang barang jadi belum dapat dieksekusi secara otomatis, sehingga ketergantungan pada tenaga manual di area kritis tersebut masih sangat tinggi.

Guna memperoleh gambaran kondisi operasional yang akurat dan kontekstual, tim TA252601001 melaksanakan wawancara daring pada tanggal 16 Oktober 2025 bersama Bapak Falih Muhtadi selaku \textit{Staff Logistics and Business Integration Systems} Paragon Corp. Dari hasil wawancara tersebut, dikonfirmasi bahwa permasalahan utama pada gudang barang jadi meliputi tiga hal: pertama, data stok tidak bersifat \textit{real-time} karena WMS belum terhubung langsung dengan sistem pusat; kedua, proses \textit{inbound} barang yang lambat dan perhitungan yang tidak presisi akibat ketergantungan pada proses manual; serta ketiga, proses \textit{picking} yang masih bergantung pada operator dan rawan \textit{human error}, khususnya dalam konteks jumlah barang. Bapak Falih juga mengonfirmasi adanya permasalahan \textit{mistracking} barang akibat area penempatan sementara yang tidak tercatat di sistem, serta potensi besar untuk mengotomasi proses \textit{inbound}, khususnya pemindahan barang dari truk ke \textit{pallet changer}. Berdasarkan seluruh temuan tersebut, tugas akhir ini mengusulkan pengembangan \textit{Warehouse Management Controller} (WMC) sebagai layanan \textit{backend} berbasis arsitektur \textit{microservices} yang mencakup \textit{Authentication Service}, \textit{Goods Detail Service}, \textit{Fleet Order Service}, dan antarmuka pengguna (\textit{User Interface}), guna menjamin konsistensi data stok secara otomatis serta mengoordinasikan armada \textit{mobile robot} dalam satu platform terpusat di PT Paragon Technology and Innovation.
